\chapter{Pubertät}
\label{chap:heim:pubertaet}

\section{Auffrischung \& Impulskontrolle}

% ── Übung: Abruf im Freien ────────────────────────────────────────────────────
\begin{uebung}{Abruf im Freien}{pubertaet}{Abruf}{3}
  \uebungsbeschreibung{%
    Der Hund wird auf einer langen Leine abgesetzt und darf schnüffeln. Nach einer kurzen Pause wird er mit dem Rufnamen und Abrufkommando zurückgerufen.\\[6pt]
    \textbf{Schritte:}
    \begin{enumerate}[leftmargin=*, itemsep=2pt]
      \item Hund absetzen und etwas Distanz aufbauen (3–5\,m).
      \item Einmal deutlich rufen: „\textit{[Name]} – Hier!".
      \item Kommando nicht wiederholen – ggf. mit Leine leicht unterstützen.
      \item Großes Fest bei jedem Kommen!
    \end{enumerate}
  }
  \uebungsvariationen{%
    \begin{itemize}[leftmargin=*, itemsep=4pt]
      \item \textbf{Rückwärtslaufen:} Trainer läuft bei Abruf rückwärts weg – motiviert den Hund mehr.
      \item \textbf{Versteckspiel:} Trainer versteckt sich hinter einem Baum bevor er ruft.
      \item \textbf{Ohne Leine} (nur in gesichertem Bereich und mit zuverlässigem Abruf).
    \end{itemize}
  }
  \uebungstrainer{%
    \begin{itemize}[leftmargin=*, itemsep=4pt]
      \item Niemals beim Kommen schimpfen – auch wenn es lange dauerte.
      \item Nie nur am Ende der Runde abrufen (Hund lernt: Abruf = Spaß ist vorbei).
      \item Bei Pubertätshunden erhöhte Ablenkung einkalkulieren – Erwartungen anpassen.
    \end{itemize}
  }
\end{uebung}

% ── Weitere Übungen hier einfügen ─────────────────────────────────────────────
